\documentclass[10pt,a4paper]{amsart}
\usepackage[margin=19mm]{geometry}
\usepackage{amsmath,amssymb,mathtools}
\usepackage[scr=rsfs,cal=euler]{mathalfa}
\usepackage{xfrac}
\usepackage[unicode,hidelinks,bookmarks=false]{hyperref}
\newcommand{\ie}{i.e.\ }
\newcommand{\cf}{cf.\ }
\newcommand{\resp}{respectively\ }
\newcommand{\Subsection}{\subsection}
\providecommand{\nobreakdash}{\nobreak\hbox{-}}
\setlength{\emergencystretch}{2em}
\multlinegap0pt
\usepackage{bm}
\usepackage{bbm}
\usepackage{dsfont}
\usepackage{xspace}
\usepackage{cancel}
\usepackage{mathtools}
\usepackage{amsthm}
\makeatletter
\@ifundefined{theorem}{\newtheorem{theorem}{Theorem}}{}
\@ifundefined{definition}{\newtheorem{definition}[theorem]{Definition}}{}
\@ifundefined{lemma}{\newtheorem{lemma}[theorem]{Lemma}}{}
\makeatother
\DeclareMathOperator{\im}{{\mathrm{Im}}}
\newcommand{\be}{\begin{equation}}
\newcommand{\cal}{\mathcal}
\newcommand{\ee}{\end{equation}}
\renewcommand{\bar}{\overline}
\title[Delocalization of One-Dimensional Random Band Matrices]{Delocalization of One-Dimensional Random Band Matrices}
\author{Horng-Tzer Yau, Jun Yin}
\date{Source: arXiv:2501.01718v4 (CC BY 4.0)}
\begin{document}
\maketitle

\noindent\emph{Source and licence.} Delocalization of One-Dimensional Random Band Matrices. Version arXiv:2501.01718v4 (CC BY 4.0), distributed under the Creative Commons Attribution 4.0 International licence, as recorded in the arXiv licence metadata for this exact version and shown on the abstract page https://arxiv.org/abs/2501.01718v4. Full rights notice: licenses/arxiv-2501.01718-CC-BY-4.0.txt.\par\smallskip

\noindent\textit{Editorial scope.} Counted unit: the Lemma whose source label is \texttt{lem\_WI\_K} in the article (source0.tex lines 2613--2637, source label \texttt{lem\_WI\_K}), delivered together with the article's own complete original proof of it (source0.tex lines 2673--2970, one \texttt{proof} environment of 298 lines). No mathematics is written, repaired or re-derived: statement and proof are the article's own text line for line. The proof is detached from the statement in the source: the article states the result at lines 2613--2637 and prints its proof at lines 2673--2970, and the proof environment's own header names the statement, so the association is the article's own. Required context retained uncounted: def\_mtzk (lines 1006--1011); Def\_Ktza (lines 1290--1316); WI\_2G (lines 2602--2605). Every definition, display and label of those lines is retained unchanged. Omitted from this package and not redistributed: 1--1005, 1012--1289, 1317--2601, 2606--2612, 2638--2672, 2971--7008. Nothing inside a retained excerpt is omitted or altered: every retained body is delivered exactly as the article's lines are written, so no omission receipt of any kind accompanies this package. Citations: the retained excerpts carry no citation command at all, so no bibliography entry of the article is transcribed and no reference is redistributed. Reference adaptation: none was needed and none was made; every reference inside the delivered unit resolves inside it, so no unresolved reference or citation is produced. Printed slips kept literal: no printed word, symbol, hypothesis or number of the counted statement or of its proof was altered. Layout and class: the article's own class is \texttt{article} with packages including tikz, graphicx, biblatex, geometry, bbm, amssymb, amsthm, mathtools, amsfonts, latexsym, amsmath, amscd, amsxtra, enumerate; the wrapper is the lane's pinned \texttt{amsart} editorial wrapper at 10pt on A4, which restates the theorem-like environments the retained text needs and adds the article's own macro definitions that the retained text needs (\texttt{\textbackslash bar}, \texttt{\textbackslash be}, \texttt{\textbackslash cal}, \texttt{\textbackslash ee}, \texttt{\textbackslash im}), with extra wrapper packages (bm, bbm, dsfont, xspace, cancel, mathtools). The delivered numbering is the unit's own. Assets: no retained span contains a figure environment, a graphics-inclusion command, a PDF-page inclusion, a table or a TikZ picture; no image file is copied into this package, the package declares no asset dependency and no image file is redistributed with it. Licence: version arXiv:2501.01718v4 of this article is distributed under the Creative Commons Attribution 4.0 International licence, as recorded in the arXiv licence metadata for this exact version and shown on the abstract page https://arxiv.org/abs/2501.01718v4. A full rights notice is delivered at licenses/arxiv-2501.01718-CC-BY-4.0.txt. Characters: every retained excerpt is pure ASCII, so the delivered engine, XeLaTeX with its default Latin Modern text font, reports no missing character.
\begin{equation}\label{def_mtzk}
m (\sigma ):=m^{(E)}(\sigma ) = \begin{cases}
     m^{(E)}, & \sigma  =+ \\
     \overline{m^{(E)}} , & \sigma = -
\end{cases},
\end{equation}

\begin{definition}[The primitive equation]\label{Def_Ktza}
For 
$$
\boldsymbol{\sigma}=(\sigma_1, \sigma_2, \cdots \sigma_n),\quad 
\boldsymbol{a}=(a_1, a_2, \cdots a_n)
,\quad \sigma_i \in \{+,-\}, \quad
a_i\in \mathbb Z_L, \quad 1\le i\le n
$$
and  $m(\sigma)$ defined in \eqref{def_mtzk}, we define  ${\cal K}_{t, \boldsymbol{\sigma}, \textbf{a}}$ to be the unique solution to the equation 
 \begin{align}\label{pro_dyncalK}
       \frac{d}{dt}\,{\cal K}_{t, \boldsymbol{\sigma}, \textbf{a}} 
       =  
       {W}\cdot  \sum_{1 \le k < l \le n} \sum_{a, b} \; \left( {\cal G}^{(a), L}_{k, l} \circ {\cal K}_{t, \boldsymbol{\sigma}, \textbf{a}} \right) S^{(B)}_{ab} \left( {\cal G}^{(b), R}_{k, l} \circ {\cal K}_{t, \boldsymbol{\sigma}, \textbf{a}} \right)  
    \end{align}
with the initial value 
$$
    {\cal K}_{0, \boldsymbol{\sigma}, \textbf{a}} =  {W^{-n+1}} \cdot \prod_{k=1}^n m (\sigma_k) \cdot \textbf{1}(a_1 = a_2 = \cdots = a_n).
$$
Here we define the ${\cal G}$ operator acting on ${\cal K}$ in the same way as it acts on ${\cal L}$,  i.e., ,  
  \be\label{calGonIND}
     {\cal G}^{(a),\,L }_{k,l}  \circ {\cal K}_{t, \boldsymbol{\sigma}, \textbf{a}} = {\cal K}_{t, \; {\cal G}^{(a),\,L }_{k,l}  (\boldsymbol{\sigma}, \textbf{a})}  \quad \forall t, k, \boldsymbol{\sigma}, \textbf{a},   
       \ee
and similarly for  ${\cal G}^{(b),\, R }_{k,l}$. 
For the special case $n=1$, we define % the r.h.s. of \eqref{pro_dyncalK} is zero. Therefore for $n=1$ case
$
{\cal K}_{t,\,+,\, a}=m,\quad {\cal K}_{t,\,-,\, a}=\overline m$ for $0\le t\le 1$. 
\end{definition}

\begin{equation}
    \label{WI_2G}
G_t(+)\cdot G_t(-)= \frac{G_t(+)-G_t(-)}{2\eta_t}=\frac{G_t(+)-G_t(-)}{ 2 (1-t)\im m},\quad \eta_t=\im z_t.
\end{equation}

\begin{lemma}\label{lem_WI_K} For $n-G$ loop ${\cal L}_{t, \boldsymbol{\sigma}, \textbf{a}}$ with 
$\sigma_1=+$ and $\sigma_n=-$, 
we have
\begin{equation}\label{WI_calL}
  \sum_{a_n}{\cal L}_{t, \boldsymbol{\sigma}, \textbf{a}}=
\frac{1}{2W\eta_t}\left({\cal L}_{t, \; \boldsymbol{\sigma}+,  \; \textbf{a}/a_n}- {\cal L}_{t, \; \boldsymbol{\sigma}-,  \; \textbf{a}/a_n}\right)  
\end{equation}
and
\begin{equation}\label{WI_calK}
  \sum_{a_n}{\cal K}_{t, \boldsymbol{\sigma}, \textbf{a}}=
\frac{1}{2W\eta_t}\left({\cal K}_{t, \; \boldsymbol{\sigma}+,  \; \textbf{a}/a_n}
- {\cal K}_{t, \; \boldsymbol{\sigma}-,  \; \textbf{a}/a_n}\right)  
\end{equation}
where 
\begin{itemize}
\item  $\boldsymbol{\sigma}\pm$ is obtained by removing $\sigma_n$ from $\boldsymbol{\sigma}$ and replacing $\sigma_1$ with $\pm$, i.e., 
    $$
\boldsymbol{\sigma}\pm=(\pm, \sigma_2, \sigma_3,\cdots \sigma_{ n-1}).$$
Notice that  the length of $\boldsymbol{\sigma}\pm$ is $n-1$. 
\item  $\textbf{a}/a_n$ is obtained by removing $a_n$ from $\textbf{a}$:
 $$
\textbf{a}/a_n=(a_1, a_2,a_3,\cdots, a_{n-1})
$$
\end{itemize}
\end{lemma}

 \begin{proof}[Proof of Lemma \ref{lem_WI_K}]
 The equation \eqref{WI_calL} for  $\cal L$ follows from \eqref{WI_2G} directly.  For the \eqref{WI_calK}, we consider first that  $n=2$. 
By the explicit formula for $\cal K$ in this case, we have 
 \begin{align}\label{hsgw2}
 \sum_{a_2} {\cal K}_{t, (+,-), (a_1,a_2)}
 =\frac1W \sum_{a_2}
\left(\frac{|m^2|}{1-t|m|^2S^{(B)}}\right)_{a_1a_2}
=\frac{1}{W(1-t)}    
 \end{align}
 On the other hand, by definition, we have 
 $$
 \frac{1}{2W\eta_t}\left({\cal K}_{t, (+ ), ( a_1)}-{\cal K}_{t, (- ), ( a_1)}\right)=\frac{m-\bar m}{2W\eta_t}=\frac{1}{W(1-t)}, \quad m:= m^{(E)}
 $$
 Here we used $|m|=1$ and $\eta_t=(1-t)\im m$ .  Combining  these two identities, we have proved   \eqref{WI_calK} for $n=2$. 
 Similarly, the case $n=3$ follows from a direct calculation and  the following identities 
 $$
 \left(\frac{|m^2|}{1-t|m|^2S^{(B)}}\right)-\left(\frac{m^2}{1-tm^2S^{(B)}}\right)=
 \left(1-m^2\right)\cdot\left(\frac{1}{1-t|m|^2S^{(B)}}\right)\cdot \left(\frac{1}{1-tm^2S^{(B)}}\right).
 $$
 and
 $$
 1-m^2=m(\bar m-m)
 $$
 
 For $n\ge 4$,   we will use the primitive equation instead of the tree representation. For simplicity, we temporally denote
 $$
{\cal K}^{(*)}_{t, \boldsymbol{\sigma}, \textbf{a}}
=\sum_{a_1}{\cal K} _{t, \boldsymbol{\sigma}, \textbf{a}},\quad
{\cal K}^{(**)}_{t, \boldsymbol{\sigma}, \textbf{a}}
=\frac1{2W\eta_t}\left({\cal K}^{(+)}_{t, \boldsymbol{\sigma}, \textbf{a}}-
{\cal K}^{(-)}_{t, \boldsymbol{\sigma}, \textbf{a}}\right),
$$
\begin{equation}\label{def_calKpm}
 {\cal K}^{(+)}_{t, \boldsymbol{\sigma}, \textbf{a}}: = {\cal K}_{t, \; \boldsymbol{\sigma}+,  \; \textbf{a}/a_n}
\quad\quad 
{\cal K}^{(-)}_{t, \boldsymbol{\sigma}, \textbf{a}}: = {\cal K}_{t, \; \boldsymbol{\sigma}-,  \; \textbf{a}/a_n}
   \end{equation}
Our goal is  to prove ${\cal K}^{(*)}={\cal K}^{(**)}$.  With
 $$
 m_n=\bar m, \quad m_1= m, \quad |m|=1, \quad \eta_t\big|_{t=0}=\im m. 
 $$
By Definition  \ref{Def_Ktza},    
 $$
  \sum_{a_n}{\cal K}_{0, \boldsymbol{\sigma}, \textbf{a}}=W^{-n+1}\cdot \left(\prod_{i=1}^n m_i \right){\bf 1}(a_1=a_3=\cdots =a_{n-1}) 
  =
  \frac{1}{2W\eta_0}\left(
  {\cal K}^{(+)}_{0,\; \boldsymbol{\sigma},\; \textbf{a}}
  - 
  {\cal K}^{(-)}_{0,\; \boldsymbol{\sigma},\; \textbf{a}}
  \right) .
 $$
This implies that 
 $${\cal K}^{(*)}_{0, \boldsymbol{\sigma}, \textbf{a}} ={\cal K}^{(**)}_{0, \boldsymbol{\sigma}, \textbf{a}}.$$
 In the remainder of this subsection,  under the inductive assumption that ${\cal K}^{(*)} ={\cal K} ^{(**)}$ holds for ${\cal K}^{(*)}$ of lengths strictly less than $n$,  we will  prove  the following identity: 
  \begin{align}\label{finK*-K}
     \frac d{dt}\left({\cal K}^{(*)}-{\cal K}^{(**)}\right)_{t, \boldsymbol{\sigma}, \textbf{a}} &\;=\;\sum_{ k=1}^{n-1} \sum_{a\;b}
        \left[\left({\cal K}^{(*)}-{\cal K}^{(**)}\right)_{t, \boldsymbol{\sigma}, \textbf{a}}\Bigg|(a_k\to a)\right]
        \cdot S^{(B)}_{ab}\cdot  {\cal K}_{t, ( \sigma _k,  \sigma _{k+1}), (a_k,b)}
        \\\nonumber
       &\;+\; \frac1{1-t}\cdot \left({\cal K}^{(*)}-{\cal K}^{(**)}\right)_{t, \boldsymbol{\sigma}, \textbf{a}} 
 \end{align}
 where $(a_k\to a)$ means replacing $a_k$ in $\textbf{a}$ by $a$.  Together with ${\cal K}^{(*)}-{\cal K}^{(**)}=0$ at $t=0$, this linear differential equation  has only the trivial solution, i.e., 
$$
{\cal K}^{(*)}_{t, \boldsymbol{\sigma}, \textbf{a}} ={\cal K}^{(**)}_{t, \boldsymbol{\sigma}, \textbf{a}}, \quad 0\le t\le 1.
$$
 
 For the l.h.s. of \eqref{finK*-K}, by the primitive equation for  $\cal K$, we have 
  \begin{align} \label{gsasawwe}
     \frac{d}{dt} {\cal K}^{(*)}_{t, \boldsymbol{\sigma}, \textbf{a}}=\sum_{a_n} \frac{d}{dt}\,{\cal K}_{t, \boldsymbol{\sigma}, \textbf{a}} 
       =  
       {W}\sum_{a_n}  \sum_{1 \le k < l \le n}  
       \sum_{a\;b} 
        {\cal J}^{L}_{k,l,a}\cdot{\cal J}^{R}_{k,l,b}\cdot S^{(B)}_{ab},  
    \end{align}
    where %we temporally denote ${\cal J}^{L}$ and ${\cal J}^{R}$ as follows
    $$
    {\cal J}^L_{k,l,a}:= \left( {\cal G}^{(a), L}_{k, l} \circ {\cal K}_{t, \boldsymbol{\sigma}, \textbf{a}} \right), \quad
    {\cal J}^R_{k,l,b}:=\left( {\cal G}^{(b), R}_{k, l} \circ {\cal K}_{t, \boldsymbol{\sigma}, \textbf{a}} \right).
    $$
By definition, the index $a_n$ appears  in above $  {\cal J}^{L} $.    On the other hand,  $\frac{d}{dt} \eta_t^{-1}=(1-t)^{-1}\eta_t^{-1}$ and thus 
  \begin{align} \label{gsasawII}
     \frac{d}{dt} {\cal K}^{(**)}_{t, \boldsymbol{\sigma}, \textbf{a}}=\frac1{1-t}\cdot {\cal K}^{(**)}_{t, \boldsymbol{\sigma}, \textbf{a}}+ 
       \frac{W}{2W\eta_t}  \sum_{1\le k<l\le n-1} 
       \sum_{a\;b} 
       \left(
       {\cal J}^{L,(+)}_{k,l,a}\cdot{\cal J}^{R,(+)}_{k,l,b}
       -
      {\cal J}^{L,(-)}_{k,l,a}\cdot{\cal J}^{R,(-)}_{k,l,b}
       \right)\cdot S^{(B)}_{ab},
       \quad   
    \end{align}
 where (recall $\cal K^{(\pm)}$ defined in \eqref{def_calKpm})
 $$
    {\cal J}^{L,(\pm)}_{k,l,a}:= 
    \left( {\cal G}^{(a), L}_{k, l} \circ {\cal K}^{(\pm)}_{t, \boldsymbol{\sigma}, \textbf{a}} \right), \quad
    {\cal J}^{R,(\pm)}_{k,l,b}:=\left( {\cal G}^{(b), R}_{k, l} \circ  {\cal K}^{(\pm)}_{t, \boldsymbol{\sigma}, \textbf{a}} \right).
$$
We now demonstrate that the right-hand sides of \eqref{gsasawwe} and \eqref{gsasawII} are identical, up to terms involving \(({\cal K}^* - {\cal K}^{**})\). Specifically, the special case where \(k = 1\) and \(l = n\) in \eqref{gsasawwe} contributes the term \(\frac{1}{1-t} \cdot {\cal K}^{(*)}_{t, \boldsymbol{\sigma}, \textbf{a}}\). In most other cases, we find that 
$
{\cal J}^{R,(+)}_{k,l,b} = {\cal J}^{R,(-)}_{k,l,b} = {\cal J}^{R}_{k,l,b},
$
while 
$
{\cal J}^{L}_{k,l,a}
$  
can be expressed using 
$
{\cal J}^{L,(\pm)}_{k,l,a}
$
by inductive assumption. For the remaining few cases (e.g., \(k = 1\), \(l = 2\)), direct cancellations occur between the right-hand sides of \eqref{gsasawwe} and \eqref{gsasawII}. The following proof provides a detailed, albeit tedious,  argument. 
%Readers interested in the key ideas may construct the proof for the cases \(n = 4\) and \(n = 5\) for clarity.
 



The following  identities  can be easily verified from their definitions. These identities rely on the fact that ${\cal K}^{(+)}$ and ${\cal K}^{(-)}$ differ only in the first component of $\sigma$  being $+$ or  $-$. Similarly, ${\cal K}^{(\pm)}$ and ${\cal K}$ differ slightly in their definitions. 

\begin{itemize}
    \item  For $k,l\in[[2, n-1]]$,  
 \begin{equation}\label{kdjuwsz}
  {\cal J}^{R,(+)}_{k,l,b}=  {\cal J}^{R,(-)}_{k,l,b}={\cal J}^{R }_{k,l,b} ,\quad k,l\in[[2, n-1]]   
 \end{equation}

\item For $k,l\in[[2, n-1]]$ and $ l-k  \ge 2$, the length of  $ {\cal J}^{L }_{k,l,a} $ is no longer than  $ n-1$. By induction, we have 
\begin{equation}\label{a';slw}
    \sum_{a_n}{\cal J}^{L}_{k,l,a}-\frac{1}{2W\eta_t}\left({\cal J}^{L,(+)}_{k,l,a} -{\cal J}^{L,(-)}_{k,l,a} \right)=0 ,\quad  k,l\in[[2, n-1]], \;  l-k \ge 2
\end{equation}


\item For $k,l\in[[2, n-1]]$ and $l=k+1$,
 $$
{\cal J}^L_{k,l,a} ={\cal K}_{t, \boldsymbol{\sigma}, \textbf{a}}\Big|(a_k\to a)
 ,\quad  {\cal J}^R_{k,l,b}={\cal K}_{t, (\sigma_k, \sigma_{k+1}), (a_k, b)},\quad k,l\in[[2, n-1]],\quad l=k+1
$$
 
Similarly, 
$$
{\cal J}^{L,(\pm)}_{k,l,a} ={\cal K}^{(\pm)}_{t, \boldsymbol{\sigma}, \textbf{a}}\Big|(a_k\to a).
% ,\quad  {\cal J}^{R,(\pm)}_{k,l,b} ={\cal K}_{t, ( \sigma _k,  \sigma _{k+1}), (a_k, b)}.
$$
Therefore, under the same conditions on $k, l$, 
\begin{equation}\label{uaiw03}
    \sum_{a_n}{\cal J}^{L}_{k,l,a}-\frac{1}{2W\eta_t}\left({\cal J}^{L,(+)}_{k,l,a} -{\cal J}^{L,(-)}_{k,l,a} \right)=
    \left({\cal K}^{(*)}-{\cal K}^{(**)}\right)_{t, \boldsymbol{\sigma}, \textbf{a}}\Bigg|(a_k\to a). %\quad \quad  k,l\in[[2, n-1]], \;  l-k=1
\end{equation}

 
\end{itemize}

Combining the identities \eqref{uaiw03}, \eqref{kdjuwsz} and \eqref{a';slw}, we can bound the following parts in \eqref{gsasawII} and \eqref{gsasawwe} with $\cal K^{(*)}-\cal K^{(**)}$ 
\begin{align}\label{WI-pr-Mid}
    & \sum_{  2\le k<l\le n-1} \sum_{a\;b}
    \left( {W}\sum_{a_n} 
        {\cal J}^{L}_{k,l,a}\cdot{\cal J}^{R}_{k,l,b}
        -\frac{W}{2W\eta_t} 
        \left(
        {\cal J}^{L,(+)}_{k,l,a}\cdot{\cal J}^{R,(+)}_{k,l,b}
        -
        {\cal J}^{L,(-)}_{k,l,a}\cdot{\cal J}^{R,(-)}_{k,l,b}
        \right)
        \right)\cdot S^{(B)}_{ab} 
        \\\nonumber
      =  & \sum_{ k=2 }^{n-2} \sum_{a\;b}
        \left[\left({\cal K}^{(*)}-{\cal K}^{(**)}\right)_{t, \boldsymbol{\sigma}, \textbf{a}}\Bigg|(a_k\to a)\right]
        \cdot S^{(B)}_{ab}\cdot  {\cal K}_{t, ( \sigma _k,  \sigma _{k+1}), (a_k,b)}
\end{align}

\bigskip

Next, we estimate the cases that $k=1$ or $l=n$. 

 
\begin{itemize}
\item For $k=1$ and $l=n$, we have 
$$
{\cal J}^L_{k,l,a}  =
 {\cal K}_{t, \,(+,\,-),\, (a,\,a_n)},\quad 
 {\cal J}^R_{k,l,b}={\cal K}_{t, \boldsymbol{\sigma}, \textbf{a}}\Bigg|(a_n\to b)
 $$
 and thus 
 \begin{align}\label{nssfku}
     {W}\sum_{a_n} 
        {\cal J}^{L}_{k,l,a}\cdot{\cal J}^{R}_{k,l,b}\cdot S_{ab}^{(B)}=\frac{1}{1-t}\cdot {\cal K}^{(*)}_{t,\boldsymbol{\sigma},\textbf{a}}.
 \end{align}
 

\item 
\begin{equation}\label{consjahz}
    k=1,\; l=m,  \quad \hbox{and}\quad   k=m, \;  l=n
\end{equation} 
for some $3\le m\le n-2 $. By induction, 
\begin{align}\label{awslw22}
    &\sum_{a_n}{\cal J}^{L}_{k,l,a}-\frac{1}{2W\eta_t}\left({\cal J}^{L,(+)}_{1,\,m,\,a} -{\cal J}^{L,(-)}_{1,\,m,\,a} \right)=0,\quad  k=1, \; l=m
    \\\nonumber
    &\sum_{a_n}{\cal J}^{L}_{k,l,a}-\frac{1}{2W\eta_t}\left({\cal J}^{R,(+)}_{1,\,m,\,a} -{\cal J}^{R,(-)}_{1,\,m,\,a} \right)=0,\quad  k=m, \; l=n
\end{align} 
and 
\begin{align}\label{awslw33}
    &{\cal J}^R_{k,l,b}={\cal J}^{R,(+)}_{1,m,b},\quad  k=1, \; l=m
    \\\nonumber
    &{\cal J}^R_{k,l,b}={\cal J}^{L,(-)}_{1,m,b},\quad  k=m, \; l=n
\end{align}
Hence  for calculating the  $\sum_{ab}\sum_{a_n} 
        {\cal J}^{L}_{k,l,a}\cdot{\cal J}^{R}_{k,l,b}S_{ab}$, one will see the following terms for $(k,l)=(1,m)$ and $(k,l)=(m,n)$, 
  \begin{align}\nonumber
   \frac{1}{2W\eta_t}   \sum_{ab}  {\cal J}^{L,(-)}_{1,\,m,\,a}\cdot {\cal J}^{R,(+)}_{1,\,m,\,b} S_{ab}\quad  & for \quad (k,l)=(1,m)\\\nonumber
\quad  \frac{-1}{2W\eta_t} \sum_{ab}  {\cal J}^{R,(+)}_{1,\,m,\,a}\cdot {\cal J}^{R,(-)}_{1,\,m,\,b} S_{ab}\quad &  for \quad (k,l)=( m,n)
  \end{align}
   They  cancel each other,  therefore %in this case, similar to \eqref{WI-pr-Mid}, we have 
 \begin{align}\label{uisiip}
    \sum_{k,l}^*  \sum_{a\;b}
    \left( {W}\sum_{a_n} 
        {\cal J}^{L}_{k,l,a}\cdot{\cal J}^{R}_{k,l,b}
        -\frac{W}{2W\eta_t} 
        \left(
        {\cal J}^{L,(+)}_{k,l,a}\cdot{\cal J}^{R,(+)}_{k,l,b}
        -
        {\cal J}^{L,(-)}_{k,l,a}\cdot{\cal J}^{R,(-)}_{k,l,b}
        \right)
        \right)\cdot S^{(B)}_{ab} =0; 
 \end{align}
 here  $\sum_{kl}^*$ denote   summing  over  $k,l$ satisfying  \eqref{consjahz}.  

 \item Since we assume that  $n\ge 4$, there are only four cases left, i.e., 
 \begin{equation}\label{keayew}
  (k,l)=(1,2), \; (1,n-1), \;(2,n), \;(n-1, n).
 \end{equation}
 
By definition,  
\begin{align} 
 (k,l)=(1,2) \quad &\quad   {\cal J}^L_{k,l,a}  ={\cal K}_{t, \boldsymbol{\sigma}, \textbf{a}}\Bigg|(a_1\to a) ,  
 &{\cal J}^R_{k,l,b}=  {\cal J}^{R,(+)}_{1,\,2,\,b}\quad 
 \\\nonumber
 (k,l)=(1,n-1)\quad &\quad {\cal J}^L_{k,l,a}  =
 {\cal K}_{t, \,(+,\;\sigma_{n-1},\,-),\, (a,\,a_{n-1},\,a_n)},  
 &{\cal J}^R_{k,l,b}=  {\cal J}^{R,(+)}_{1,\,{n-1},\,b}
  \\\nonumber
 (k,l)=(2,n)  \quad &\quad  {\cal J}^L_{k,l,a}  =
 {\cal K}_{t, \,(+,\;\sigma_{2},\,-),\, (a_1,\,a,\,a_n)} ,  
 &{\cal J}^R_{k,l,b}=  {\cal J}^{L,(-)}_{1,\,2,\,b}\quad 
 \\\nonumber
 (k,l)=(n-1,n)\quad &\quad {\cal J}^L_{k,l,a}  ={\cal K}_{t, \boldsymbol{\sigma}, \textbf{a}}\Bigg|(a_{n-1}\to a), 
 &{\cal J}^R_{k,l,b}=  {\cal J}^{R,(+)}_{1,\,{n-1},\,b}
\end{align}
 Summing  up $a_n$ and multiplying  $W$,  we obtain that
 \begin{align} 
 (k,l)=(1,2) \quad &\quad  W\sum_{a_n} {\cal J}^L_{k,l,a}  ={\cal K}^{(*)}_{t, \boldsymbol{\sigma}, \textbf{a}}\Bigg|(a_1\to a) ,  
   \\\nonumber
 (k,l)=(1,n-1)\quad &\quad W \sum_{a_n}{\cal J}^L_{k,l,a}  =
 \frac{W}{2W\eta_t} 
 \left({\cal K}_{t, \,(+,\;\sigma_{n-1}),\, (a,\,a_{n-1} )} 
 - {\cal K}_{t, \,( \sigma_{n-1},-),\, ( a_{n-1}, \, a )}  \right)
 \\\nonumber
 & \quad \quad \quad\quad \quad \quad\;=  \frac{W}{2W\eta_t} 
 \left({\cal J}^{L, (+)}_{1,n-1,a} 
 - {\cal J}^{L, (-)}_{1,n-1,a} \right)
   \\\nonumber
 (k,l)=(2,n)  \quad &\quad  W\sum_{a_n} {\cal J}^L_{k,l,a}  =
\frac{W}{2W\eta_t} 
 \left({\cal K}_{t, \,(+,\;\sigma_{2}),\, (a,\,a_{2} )} 
 - {\cal K}_{t, \,( \sigma_{2},-),\, ( a_{2}, \, a )}  \right)
 \\\nonumber
 & \quad \quad \quad\quad \quad \quad\;=  \frac{W}{2W\eta_t} 
 \left({\cal J}^{R, (+)}_{1,\,2,\,a} 
 - {\cal J}^{R, (-)}_{1,\,2,\,a} \right)  \\\nonumber
 (k,l)=(n-1,n)\quad &\quad  W\sum_{a_n}{\cal J}^L_{k,l,a}  ={\cal K}^{(*)}_{t, \boldsymbol{\sigma}, \textbf{a}}\Bigg|(a_{n-1}\to a), 
 \end{align}
  On the other hand, we have 
  \begin{align}
      {\cal K}^{(**)}_{t, \boldsymbol{\sigma}, \textbf{a}}\Bigg|(a_1\to a) 
      & = \frac{W}{2W\eta_t}
      \left({\cal J}^{L, (+)}_{1,\,2,\,a}-{\cal J}^{L, (-)}_{1,\,2,\,a}\right)
      \\\nonumber
       {\cal K}^{(**)}_{t, \boldsymbol{\sigma}, \textbf{a}}\Bigg|(a_{n-1}\to a) 
      & = \frac{W}{2W\eta_t}
      \left({\cal J}^{R, (+)}_{1,\,n-1,\,a}-{\cal J}^{R, (-)}_{1,\,n-1,\,a}\right)
  \end{align}
  Therefore, with $\sum_{kl}^{**}$ denoting  summing  $k,l$ in \eqref{keayew}, we have  
  \begin{align}\label{uisiip222}
    &\sum_{k,l}^{**}  \sum_{a\;b}
    \left( {W}\sum_{a_n} 
        {\cal J}^{L}_{k,l,a}\cdot{\cal J}^{R}_{k,l,b}
        -\frac{W}{2W\eta_t} 
        \left(
        {\cal J}^{L,(+)}_{k,l,a}\cdot{\cal J}^{R,(+)}_{k,l,b}
        -
        {\cal J}^{L,(-)}_{k,l,a}\cdot{\cal J}^{R,(-)}_{k,l,b}
        \right)
        \right)\cdot S^{(B)}_{ab} 
        \\\nonumber
       &=\sum_{k}\left({\bf 1}_{k=1}+{\bf 1}_{k=n-1}\right) \sum_{a\;b}
        \left[\left({\cal K}^{(*)}-{\cal K}^{(**)}\right)_{t, \boldsymbol{\sigma}, \textbf{a}}\Bigg|(a_k\to a)\right]
        \cdot S^{(B)}_{ab}\cdot  {\cal K}_{t, ( \sigma _k,  \sigma _{k+1})_{a_k,b}}
 \end{align}
 
\end{itemize}
At last,  combining  the identities \eqref{uisiip222}, \eqref{uisiip}, \eqref{nssfku}, \eqref{WI-pr-Mid}, \eqref{gsasawwe} and \eqref{gsasawII}, we obtain  the desired result  \eqref{finK*-K} and prove the Lemma \ref{lem_WI_K} by induction. 
\end{proof}

\end{document}
